\section{Conclusions}
\label{sec:conclusions}

This study used DEM simulations of undrained cyclic torsional shear in hollow cylindrical specimens to investigate how initial stress anisotropy (\Kzero{}) affects liquefaction resistance under HCA-type boundary conditions. The main findings are summarized as follows:

\begin{enumerate}
\item An analytical solution was derived for the combined servo mechanism of the DEM-HCA model by explicitly accounting for the coupling among radial, axial, and torsional boundary controls through the constant-volume constraint and the shared dependence of stress quantities on multiple wall movements. The resulting closed-form coupled stiffness matrix determines all servo coefficients directly from aggregate contact stiffnesses at each timestep, removing the independent-control approximation adopted in earlier formulations and ensuring stable convergence under general loading conditions in undrained cyclic shear.

\item Liquefaction resistance exhibits a monotonic directional dependency on initial stress anisotropy: compression-state anisotropy (\Kzero{} $< 1.0$) produces higher resistance than the isotropic state, while extension-state anisotropy (\Kzero{} $> 1.0$) reduces it. This ordering is consistent across all CSR levels and both relative densities examined.

\item The difference in liquefaction resistance cannot be explained solely by the total cumulative shear work at liquefaction. While normalized energy accumulation patterns are similar among \Kzero{} states, the per-cycle energy input rate differs because of distinct stiffness degradation trajectories: the compression-state specimen degrades stiffness most slowly, and the extension-state specimen fastest.

\item The two microscopic descriptors play complementary roles. The mechanical coordination number \Zm{}, which is strongly coupled to macroscopic relative density, governs the baseline resistance level across density groups. The signed fabric anisotropy indicator $\alpha$ captures the directional fabric effect that differentiates \Kzero{} states within a given density: at comparable $Z_{m0}$, specimens with larger $\alpha_0$ (compression-state consolidation) consistently exhibit higher liquefaction resistance. Decomposition of the fabric tensor in the cylindrical coordinate system reveals that contacts whose normals lie on the $z$--$\theta$ shear plane directly resist the applied torsional shear $\tau_{z\theta}$; compression-state consolidation concentrates a larger fraction of contacts onto this plane, providing a direct microstructural explanation for the $\alpha_0$--$N_L$ correlation.

\item Multiple indicators, including shear wave velocity, mechanical coordination number, fabric anisotropy indicator, contact-density distributions, and force-chain configurations, all show that initially distinct fabric states progressively degrade and converge toward similar configurations in the post-liquefaction stage. This convergence indicates that cyclic loading drives the granular assembly into a common dynamic equilibrium in which the fabric oscillates with loading reversals, regardless of the initial stress anisotropy.

\end{enumerate}
